% ============================================================
%  Harald Høffding, "Religion og Videnskab"
%  Religionshistoriske Smaaskrifter I
%  København: Gyldendalske Boghandel — Nordisk Forlag, 1910
%  TRANSCRIPTION (Danish)
%  Source scan: ~/bibliotek/Høffding, Harald/religion-og-videnskab.pdf
%  Page-offset: PDF = printed + 9  (printed p.1 = PDF p.10)
%  Editorial policy:
%    - Quotation marks reproduced as printed: Danish low-high „…“
%    - Original printer's typos reproduced verbatim and flagged with % [sic]
%  Progress: batch 1 = printed pp. 1–10 (PDF 10–19). Continues mid-paragraph.
% ============================================================
\documentclass[12pt,a4paper]{article}
\usepackage[utf8]{inputenc}
\usepackage[T1]{fontenc}
\usepackage{libertinus}
\usepackage{libertinust1math}
\usepackage[danish]{babel}
\usepackage[protrusion=true,expansion=true]{microtype}
\usepackage{setspace}
\usepackage[margin=1.2in]{geometry}
\usepackage{fancyhdr}
\usepackage{hyperref}

\hypersetup{
  pdftitle={Religion og Videnskab},
  pdfauthor={Harald Høffding},
  hidelinks
}

\pagestyle{fancy}
\fancyhf{}
\rhead{Høffding, \emph{Religion og Videnskab} (1910)}
\cfoot{\thepage}

\setstretch{1.3}

\title{\textbf{Religion og Videnskab}\\[0.4em]
\large \emph{Religionshistoriske Smaaskrifter} I}
\author{Harald Høffding}
\date{København: Gyldendalske Boghandel --- Nordisk Forlag, 1910}

\begin{document}
\maketitle
\thispagestyle{fancy}

\bigskip

\noindent Indenfor enhver Kultur, vi kende, indtræder der en kritisk
Tilstand, naar Arbejdsdelingen træder i Kraft. Paa det materielle
Omraade bliver det da ikke længere muligt for den enkelte Arbejder at
frembringe et Værk, ved hvilket han kan have Følelsen af at have virket
med alle sine Evner. Hvad den Enkelte kan gøre, bliver nu oftest kun at
frembringe en enkelt Del af et Hele, en Del, der kun kræver Anvendelsen
af en enkelt Evne. Udviklingen af hans Kraft og hans Snille bliver
derfor ensidig. Som hans Værk bliver selve hans Personlighed let kun et
Brudstykke. Paa det aandelige Omraade betyder Arbejdsdelingen, at Kunst,
Videnskab, Religion og Moral, der i tidligere Tider udgjorde en Enhed,
nu træde ud og ofte træde i Modsætning til hverandre som stridende
Magter. I tidligere Kulturperioder stod Religionen som den store
centrale Livsstrøm, til hvilken alle aandelige Kræfter og Drifter gav
deres Bidrag, og som paa den anden Side ydede fuld Tilfredsstillelse for
den kunstneriske og den intellektuelle Trang, forsaavidt saadan Trang
overhovedet rørte sig. Men efterhaanden søge de aandelige Fornødenheder
deres Tilfredsstillelse ad forskellige Veje. Forskertrangen
tilfredsstilles ikke længere ved de Svar, religiøse Ideer kunne give paa
de mangehaande Spørgsmaal, Livet og Erfaringen kunne føre til at
opkaste. Der udvikles særlige Metoder, særlige Former for at stille
Spørgsmaalene og at søge Svarene. Den kunstneriske Trang gaar ligeledes
mere og mere sine egne Veje, og der søges et selvstændigt Grundlag for
moralsk Vurdering. Historien viser, at i Modsætning til disse
Nydannelser har Religionen en Tendens til at holde fast ved de en Gang
fæstnede Forestillinger og Handlemaader. Den modsætter sig den aandelige
Arbejdsdeling, saalænge den kan, og den søger saalænge som muligt at
lukke Øjnene for det aandelige Livs forandrede Vilkaar. Deraf følger en
Delthed og en Splittelse i det aandelige Liv. Der mangler den Samling og
Harmoni, den Koncentration af alle Sindets Kræfter i én fælles Retning,
der under de tidligere Forhold var forholdsvis let at opnaa. Herpaa beror
det, at der existerer et religiøst Problem, og at et saadant vil vedblive
at bestaa, indtil Aandslivet hos de Enkelte og i Slægten som Helhed har
vundet nye samlede og samlende Former til Fyldestgørelse af sin inderste
Trang, --- om saadanne nye Former overhovedet ere mulige.

Jeg vil her kun fremdrage nogle Træk af Problemets intellektuelle Side.
Det hele Spørgsmaal er for mig ikke udtømt ved at ses fra denne Side. Men
det har sin store Betydning at blive klar over, hvorpaa den Modsætning,
der atter og atter gør sig gældende mellem Religion og Videnskab, egentlig
beror. Ofte finder man --- og det gælder begge de Parter, der her staa
overfor hinanden --- Grunden til denne Modsætning paa urette Sted. Der er
megen overfladisk og dogmatisk Pukken paa Videnskabens virkelige eller
formentlige Resultater paa den ene Side, og der er paa den anden Side
megen Kortsynethed med Hensyn til den Ændring, Tankelivet faktisk er
undergaaet, ganske afset fra de mange iøjnefaldende Resultater. I de
Sammenstød, hvortil Modsætningen fører, savner Tilskueren ikke Ærlighed i
Bekendelserne fra den ene eller den anden Side. Der er jo i vore Dage
ogsaa forholdsvis liden Fare ved at aflægge sin Bekekendelse, i hvilken
% [sic] "Bekekendelse" for "Bekendelse" — printer's error in the 1910 print
Retning denne saa gaar. Derimod har man rigelig Lejlighed til at savne den
intellektuelle Redelighed, der ikke lader Sindet hvile, før de store
Spørgsmaal ere undersøgte fra alle Sider og ved alle Midler, der staa til
Raadighed. Det er desværre ikke en Dyd, den religiøse Etik særligt har
indskærpet, og mange af Religionens Modstandere vise ved Arten af deres
Polemik, at de heller ikke kende de strenge Fordringer, den stiller. Det
Savn, en Tilskuer her kan føle, hænger sammen med, at den religiøse Debat
i vore Dage mest har Agitationens Karakter. Det gælder paa begge Sider om
at vinde Tilhængere, om Udbredelse, ikke om Uddybelse. Nye Tanker har
denne moderne Debat da heller ikke bragt.

Hvad jeg her vil bringe, er heller ikke nye Tanker, men gamle Tanker, man
ikke har ladet komme til deres fulde Ret.

\begin{center}---\end{center}

Gaa vi tilbage til Tider, i hvilke Religionen ikke blot var Hjælp og
Trøst, men ogsaa gav Svar paa, hvad Tanken da kunde spørge om, --- Tider
altsaa, i hvilke den hele Verdensopfattelse havde en religiøs Karakter,
--- saa er der to Ejendommeligheder, der især fremtræde. Den hele
Verdensopfattelse har en anskuelig Karakter. Tilværelsen staar som et
Hele, hvis enkelte Regioner kunne ses eller dog forestilles, uden at det
var nødvendigt at bryde med, hvad Sanserne umiddelbart syntes at lære.
Verdensrummet, som det fremstillede sig for Sansning og Fantasi, indeholdt
Mulighed for at lokalisere alle de store Begivenheder, om hvilke den
religiøse Forkyndelse handlede. Og naar der i denne Anskuelighedens Verden
indtraadte Forandringer, især betydningsfulde Forandringer, var det ikke
ad vanskelig Forskens Vej \emph{Forklaringen} skulde søges. Saadanne
Forandringer forklaredes ved personlige Væseners Indgriben i
Begivenhedernes Gang. Det er Anskueligheden, og det er de personlige
Aarsagers Princip (hvad man i videste Forstand kan kalde Animisme), der
karakteriserer en specifik religiøs Verdensopfattelse.

Men baade Rummets Begreb og Aarsagsbegrebet have under den selvstændige
Erkendelsestrangs Indflydelse undergaaet Forandringer, der her ere af
afgørende Betydning.

\medskip

1. Vi have en naturlig Trang til at henføre Alt, hvad der sker, til
bestemte Steder i Rummet, og Religionen fornegter i sine klassiske Tider
ikke denne Trang. De religiøse Forestillinger have Anskuelsens, ikke
Begrebets Karakter, og deres store Indflydelse, deres Evne til at gribe og
fylde, beror for en ikke ringe Del derpaa. De højere Folkereligioners
Indhold er en Række af historiske Begivenheder, i hvilke Menneskenes og
Menneskelivets Skæbne er indflettet, --- et stort Drama, og et Drama, der
spilles paa en Skueplads, Sanseanskuelsen, eller i hvert Tilfælde
Fantasien meget vel kan omfatte uden at sprænges. Hele Tilværelsen,
saaledes som den paa dette Trin er tilgængelig, tages i Beslag til denne
Skueplads, og dens forskellige Regioner afgive Pladser til de forskellige
Scener i Dramet. Jorden staar som Midtpunkt eller som det laveste Sted, og
fra den synlige Himmel, der er Guddommens Verden, stige guddommelige
Væsener ned for at gribe ind i Menneskenes Skæbne, maaske for selv at lide
og dø, inden de stige op igen i de himmelske Regioner, til hvilke de
derpaa aabne Adgang for dem, der ville følge dem. „Alle Mennesker,“ siger
Aristoteles, „antage, at der er Guder, og alle, baade Hellener og
Barbarer, tildele de Guddommen det højeste Sted.“ Denne Udtalelse passer
paa alle Folkereligioner, saa forskellige de ellers kunne være. De
religiøse Forestillinger vævedes paa naturlig Maade ind i Menneskenes hele
uvilkaarligt dannede Opfattelse, og der existerede endnu ingen Videnskab,
der nødvendiggjorde en radikal Ændring af Opfattelsen. Der var ingen
Vaghed eller Utydelighed i Forestillingen om det guddommelige; dette kunde
lokaliseres ligesom alt Andet; man kunde pege mod det Sted, hvorfra de
guddommelige Kræfter fremvældede, saavel som paa det Sted, paa hvilket de
virkede. De guddommelige Tanker vare ganske vist ligesaa højt hævede over
Menneskets Tanker, som Himmelen over Jorden; men Himmelen var netop ikke
saa højt hævet over Jorden, at Anskuelsen maatte give tabt. Trods
Forskellen i Magt og Værdighed var dog et anskueligt og nært Forhold til
Guderne muligt.

Ingen Folkereligion har forladt dette Verdensbillede. Men før endnu
Videnskaben for Alvor havde rokket ved Verdens Anskuelighed, kom der dog
hos tænkende Aander Betænkelighed frem, og der rokkedes ved den Tryghed,
hvormed man hidtil havde dvælet ved det anskuelige Verdensbillede som
Rammen for alt Religionens Indhold. Og denne Uro skyldtes for en
væsentlig Del en Uddybelse af selve den religiøse Bevidsthed. Der rører
sig i al Religion en dobbelt Trang, dels en Trang til at hæve de religiøse
Forestillingers Genstand saa højt over menneskelige og jordiske Kaar og
Former som muligt, dels netop en Trang til at føle sig Et med Guddommen,
eller dog at være den nær. Det er en Trang til at staa overfor et Ophøjet,
og en Trang til et inderligt Forhold til det Højeste. Ad begge Veje kunde
der affødes et Brud med Folkereligionens naive Anskuelighed. Naar
Guddommen skal være ophøjet over Alt, maa den saa ikke ogsaa være ophøjet
over alle rumlige Forskelle, over Muligheden af Lokalisation og
Begrænsning? Og naar den skal staa i det inderligste Forhold til
Mennesket, kan den saa være adskilt i Rummet fra det Menneske, med hvilket
den skal være Et? En Afstand i Rummet er altid et udvortes Forhold, hvor
hurtigt den saa gennemløbes. I selve de oprindelige religiøse Dokumenter
spores begge Tendenser, kundgør der sig en rent religiøs Trang, der maatte
føre ud over Anskueligheden. Dermed forenede sig saa ogsaa Indflydelse af
den filosofiske Reflexion, for Kristendommens Vedkommende især
Indflydelsen af Platonismen. Det var under Indvirkning af disse
forskellige Motiver, at Augustinus kunde sige om sin Gud, at han var
„overalt og intetsteds,“ og at han havde sin Plads „længere inde i Sjælen
end dennes inderste Væsen,“ saavel som „højere oppe end Sjælen forstod at
stige.“ Men forinden var der først en haard --- baade legemlig og
aandelig --- Kamp mellem de religiøse Materialister (Antropomorfister),
der holdt paa den absolute Anskuelighed, og de religiøse Spiritualister,
der negtede denne.

Det var i Teorien Spiritualismen, der sejrede. Men man kunde og vilde ikke
helt fornegte Anskueligheden, der jo ogsaa fremtræder saa klart og
utvetydigt i de oprindelige religiøse Dokumenter og tillige er af saa stor
pædagogisk Betydning. Ved en Fortælling som den om Jesu Himmelfart er det
jo aldeles nødvendigt, naar den skal tages bogstaveligt, at opfatte Himlen
som et bestemt Sted, op til hvilket Farten gik. Himmelfarten er en
Overgang fra den jordiske Del af Rummet til den himmelske. Middelalderens
Teologer stillede sig da her saaledes, at de lærte, at Himmelfarten betød
en Overgang til en anden Art af Existens, en aandelig Tilværelse. At stige
op til Himmelen vil sige at aflægge den jordiske og menneskelige
Begrænsning. Men denne symboliserende Opfattelse udelukkede for dem ikke
den bogstavelige. Man fastholder paa en Gang Anskueligheden og
Symbolismen. Der er virkeligt sket en Overgang fra fysisk Jord til fysisk
Himmel, men selve denne Overgang betyder tillige noget Andet. Den
bogstavelige og den billedlige Betydning af „det Højere“ i Modsætning til
„det Lavere“ fastholdes begge to --- skønt man øjensynligt er tilbøjelig
til at lægge mest Vægt paa den billedlige. Man antog, at den dobbelte
Mening --- den realistiske og den idealistiske, den bogstavelige og den
opbyggelige --- havde været nedlagt i selve de oprindelige Beretninger.
Hvad der for den oprindelige religiøse Bevidsthed stod som en udelt
Sandhed, anskuelig og betydningsfuld paa én Gang og med et Slag, det stod
for den middelalderlige Bevidsthed som et Baade---Og, paa engang sondret
og uadskilleligt.

I Dante's store Digt raader dette Standpunkt. Den symbolske Betydning er
vel ogsaa her den, der ligger Forfatterens Hjerte nærmest. Men den ideelle
Sandhed har en tilsvarende ydre Dragt, hvis Realitet fastholdes. Det
Enestaaende i Dantes Divina comedia beror, fra dette Synspunkt set, netop
derpaa, at de ydre Rammer staa fast, have en Massivitet og en
Anskuelighed, som senere Tiders religiøse Poesi ikke kunde naa, og at dog
Inderligheden fuldt ud kommer til sin Ret ved den symbolske Tydning, der
skinner igennem. Overgangen fra det Materielle til det Ideelle gøres ---
naivt og stort paa en Gang --- ved selve Rummets Grænse: den yderste
Himmel, der omslutter Alt, har selv sit „Hvor?“ i Guddommen, hvis Kraft og
Kærlighed er Kilden til Verdenernes Omsving.

Men der kom den Tid, da man fordrede et mere prosaisk Svar paa
Spørgsmaalet „Hvor?“ Ligesom vi indenfor vor Erfaring bestemme ethvert
Sted ved dets Forhold til et andet Sted, saaledes maa vi ogsaa antage, at
overalt i Verden, hvor vi tænke os hen, gælder det Samme. Der kan da ingen
absolute Steder eller absolute Grænser tænkes. Ethvert anskueligt Emne har
sit „Hvor“ i bogstavelig Betydning, bestemt ved Forhold til andre Steder.
Der kan derfor intet Verdensmidtpunkt og ingen Verdensgrænse forestilles.
Der kan ikke findes noget Sted, ud fra hvilket de guddommelige Kræfter
kunde strømme ud eller ned. Paa ethvert Sted maa den samme Lov, eller
rettere den samme Metode for Stedsbestemmelse finde Anvendelse. Det var
netop denne Tankegang, der fik Giordano Bruno til at drage sine dristige
Konsekvenser af Kopernikus' Hypotese. Allerede i Lucretius' berømte
Læredigt antydedes denne Tankegang, men først langt senere blev det muligt
at holde den fast og anvende den. Overfor den uoverskuelige Rummets
Verden, der nu aabnede sig, maatte den antike Modsætning mellem Himmel og
Jord staa som forsvindende. At Rummet faktisk er uendeligt, kan i en vis
Forstand ikke bevises. Det Væsentlige er den nye Metode, den nye Maade at
tænke paa, naar Talen er om Sted og Rum. Det er paa en Maade et nyt Organ,
der udvikler sig, og som Tanken bringer med, hvor i Rummets Verden den saa
hensætter sig. Overfor denne nye Maade at tænke paa var det gamle
Verdensbillede med sine absolute Steder uholdbart. Den Lidenskab, med
hvilken det saa længe blev holdt fast, og med hvilken man fór frem mod det
nye Verdensbilledes Forkæmpere, viser da ogsaa, at man havde en Anelse om
hvad det gjaldt. Det gjaldt nemlig om intet Mindre end Opgivelse af de
religiøse Forestillingers bogstavelige Betydning. Alle Folkereligioner
havde lagt det gamle, absolute og begrænsede Verdensbillede til Grund. Det
afgav Grundlaget og Støttepunktet for den religiøse Fantasi og for den
Hjemlighedsfølelse, med hvilken Menneskeaanden kunde betragte en Verden, i
hvilken dens højeste Idealer kunde have deres „naturlige Steder“ ligesaa
vel som de ringeste materielle Elementer havde deres.

Middelalderens Baade---Og var nu ikke længere tilstrækkeligt. De
Begivenheder, om hvilke de bibelske Skrifter fortælle saa tilforladeligt
og anskueligt, maa nu paa afgørende Punkter forstaas billedligt. Man kan
ikke længere hjælpe sig med at antage baade en real og en ideal Betydning.
Fortællingen om Himmelfarten er her det typiske Exempel, fordi der i den
klart udtrykkes, at hele Universet berøres af hvad der skete i Palæstina.
En himmelsk Magt var stegen ned til Jorden og forlod den nu igen for at
vende tilbage til sit naturlige Sted.

Kun meget sjeldent synes den religiøse Bevidsthed at sætte sig levende ind
i de Spørgsmaal, der herved opstaa, og som baade ere af Betydning for den
historiske Opfattelse af de bibelske Fortællinger og for Karakteristiken
af den religiøse Verdensanskuelse. Kan Fortællingen om Himmelfarten ikke
være bogstaveligt sand, da staar den som en Legende, der kan have sin
Skønhed og sin ideelle Betydning, naar vi sætte os tilbage i Tidens
Tankegang, men som ikke længere kan optræde med dogmatisk Vægt. Man har
derved paa et væsentligt Punkt anerkendt den historisk-kritiske Metodes
Berettigelse. Ti denne Metode hviler paa den samme Forudsætning, paa
hvilken Bruno's Udstrækning af Stedsbestemmelsens Metode fra Jorden til
Himlen hvilede, den nemlig, at vi ved Undersøgelsen af det Fjerne ikke
kunne Andet end anvende den Fremgangsmaade, som vi have prøvet overfor det
Nære. Vi maa møde overfor de gamle Beretninger med den Verdensopfattelse,
vi nu have, selvom den er nok saa forskellig fra den, der ligger til Grund
for hine Beretninger. Men det er ikke blot for den historiske Metode, at
Konsekvenserne blive af Betydning. Hele den Betydning, der kan tillægges
religiøse Forestillinger maa derved blive en anden. Intet Under, at man
foretrækker at lade Sagen staa hen! Og alligevel blive selv ortodoxe
Teologer her uvilkaarligt Symbolister og Rationalister, saasnart de for
Alvor tænke sig ind i Emnet. Martensen f. Ex. fastholder bestemt, at
Beskrivelsen af Himmelfarten er at forstaa bogstaveligt. Men han mener, at
da Himmelen ikke er noget bestemt Sted, og da Meningen med Fortællingen
blot kan være den, at Jesus nu forlod den rumlige, sanselige Verden, saa
maa man antage, at Jesus kun har hævet sig et Stykke op i Luften, hvilket
skulde tjene Apostlene til et Tegn paa, at han nu antog en helt anden Art
af Existens end den sanselige! Til denne smagløse Tydning af den gamle
Beretning ledes Martensen ved sin Paastand om, at Bibelen ved Himlen ikke
mener et Sted, til hvilket man kommer ad lange Veje. Han gaar ud fra den
symbolske Betydning af Forestillingen om Himlen --- og dog vil han tage
den gamle Fortælling bogstaveligt. Hvad der var et muligt Standpunkt for
en middelalderlig Teolog, det var et umuligt Standpunkt i det nittende
Aarhundrede, eller burde have været det.

Der er intet Andet at gøre end at lægge hele Vægten paa den symbolske
Betydning. Det er den eneste Maade, paa hvilken de gamle Beretninger kunne
komme til deres fulde Ret, baade religiøst, poetisk og historisk. I enhver
„opbyggelig“ Anvendelse gaar man i Virkeligheden denne Vej. Kun saaledes
kan man ret forstaa hvad der har bevæget sig i de Menneskesind, hvem disse
Beretninger skyldes. De „naturlige Steder“ for de religiøse Begivenheder
ere at finde i Sjælenes Indre, og kun der. Alle Legender og Myter ere kun
Vidnesbyrd om, hvad der har grebet Sindene i en vis Tidsalder og sat
Fantasien i Bevægelse. Og det er den historiske Forskens Opgave at
udfinde, hvilke ydre Begivenheder der have været medvirkende ved disse
aandelige Tildragelser, og derved muliggøre en indgaaende psykologisk
Forstaaelse af dem. Selve den historiske Forsken maa igen, som allerede
bemærket, gaa ud fra de Kendemærker paa Virkelighed, som vi overalt
benytte, hvor vi søge at orientere os i Omverdenen. Hvor langt man ad
denne Vej kan komme, og hvilken religiøs Betydning Resultaterne af den
historiske Undersøgelse af Religionens Oprindelse og Udvikling kunne faa,
er et Spørgsmaal for sig, som jeg her ikke skal gaa ind paa. Personligt
tror jeg, at den saakaldte liberale Teologi i denne Henseende nærer for
store Forventninger. --- Jeg vender tilbage til Spørgsmaalet om
Indflydelsen af Tanken om Tilværelsens Uoverskuelighed paa den religiøse
Bevidsthed.

Det forholder sig ingenlunde saaledes, at det Problem, der her rejser sig,
kun skulde gælde for Folkereligionerne. Enhver Livsanskuelse, hvad enten
den er religiøs i snevrere Forstand eller ikke, maa berøres af det. Enhver
menneskelig Besindelse over Livets Værd og Livets Kaar maa møde den
Betragtning, at det menneskelige Liv med dets Interesser og Idealer, dets
Sejre og dets Nederlag, kun indtager en forsvindende Plads i den Verden, vi
kende. Vi kunne ikke længere gøre hele Verden til Skueplads for Livets
Drama. Den Skueplads, paa hvilken dette Drama udfolder sig, er begrænset og
forsvindende --- og hvilken Betydning og Værd har da hele det øvrige,
uoverskuelige Univers? Det er forunderligt at se, hvor liden Vægt, man i
Regelen har lagt paa dette Spørgsmaal, baade i Filosofien og i Teologien,
især efterat de første Indtryk af den nye Opfattelse havde fortonet sig.

For Bruno var det en aandelig Befrielse, da de snevre Verdensgrænser faldt
for hans Eftertanke. For hans poetisk-filosofiske Opfattelse stod nu
Tilværelsen som en stor Enhed, besjælet overalt af den samme Aand og baaren
af de samme Kræfter. Og paa lignende Maade opfattede ogsaa Jacob Bøhme
Sagen, da han havde arbejdet sig ud af den svare Anfegtelse, i hvilken han
stedtes, da Naturens Uendelighed først var gaaet op for ham i Modsætning
til den menneskelige Verdens Begrænsning. Endnu for Holberg stod det
saaledes, at Astronomien nu havde gjort et højere Gudsbegreb muligt end
det, ikke blot den menige Mand, men ogsaa mange Teologer have. Men senere
hen drøftes Spørgsmaalet sjeldent, og det skydes helst tilside. Positivister
og Romantikere stille sig her i Grunden paa lignende Maade som Teologerne.
Comte's Humanitetsreligion indførte en Symbolik, i hvilken det store Univers
stod som tjenende Baggrund for den Klode, paa hvilken Menneskelivet
udfoldede sig, og selv rent videnskabeligt afviste han al nærmere Forsken om
de fjernere Dele af Universet. Hegel (og hos os Heiberg) spottede over den
„slette“ Uendelighed, Udstrækningen i Rummet frembød. Man har efter Hegels
Mening overvurderet Betydningen af det kopernikanske System. Skønt hans
Spekulation skulde give os Grundformerne for \emph{Verdens}aandens Udvikling
hævdede han dog, at Jorden var den fuldkomneste Klode, fordi det var der, et
aandeligt Liv fandtes. Og Grundtvig fandt, at der var ingen Interesse ved
det uendelige Rum og det kopernikanske System, naar man ikke gik ud fra den
Antagelse, at Sol, Maane og alle Stjerner ere til for Menneskenes Skyld.

Der er her paa alle Standpunkter en Modsætning mellem Verdensanskuelse og
Livsanskuelse, der vel kan opfordre til Eftertanke. Inderligst knytter
enhver alvorlig Livsanskuelse sig til den snevre Kreds, indenfor hvilken
menneskelige Opgaver og Idealer gælde; men det er umuligt andet, end at
Blikket atter og atter maa glide ud til den uendelige Baggrund. Vi kunne
fastholde, at der i den forsvindende Del af Tilværelsen, vi bebo, har
udviklet sig Kræfter og Livsformer af ædlere Art end der kan spores i det
øvrige Univers, men vi raade ikke over Tankeformer, der kunde lade disse
Kræfter fremtræde som dem, der bestemme hele det uendelige Verdensløb. Ud
fra det, der for os har den højeste Værdi, kunne vi ikke udlede eller
forstaa Klodernes Baner og Verdensprocessernes Bølgegang. Der fremtræder
for os en Kløft mellem vore Idealers Rige og Materiens Verden, en Kløft,
hvorover der ikke kan bygges Bro. Man kan lukke Øjnene for denne
Kendsgerning, men mandigere er det at se den under Øjne og forme sin Tro og
sit Haab saaledes, at de kunne staa deres Prøve paa Virkelighedens
Kampplads. Vi maa have Mod til at fastholde, at den Strøm af Kraft og Liv,
der i sit Løb har bragt os Alt, hvad vi kende af Stort og Skønt, ogsaa
fremdeles vil formaa at bane sig sin Vej gennem Virkelighedens haarde
Masser, selvom vi ikke kunne forbinde nogen Mening med, at disse Masser kun
skulde være til for dens Skyld. ---

Den Anskuelighed, som er alle Folkereligioners Styrke, raader ikke blot med
Hensyn til Rummet, men ogsaa med Hensyn til Tiden, og paa ganske
tilsvarende Maade. Dog er det ikke i alle Folkereligioner, at
Tidsforestillingen og Forestillingen om en historisk Udvikling spiller en
væsentlig Rolle. Hvor disse Forestillinger komme frem, tyder det paa et
højere etisk Udviklingstrin. Det er nu ikke længere blot de enkelte
Individers eller de smaa Menneskegruppers øjeblikkelige og sanselige
Velfærd, der er den ledende Interesse ved religiøse Forestillingers
Udformning; det er et helt Folks, og tilsidst den hele Menneskeslægts
Frelse, det gælder. De religiøse Problemer angaa da Meningen med den hele
menneskelige Historie her paa Jorden. Religionshistorien viser os her en
ejendommelig Modsætning. Ti paa det Punkt i Udviklingen, hvor Trangen til at
finde Mening i Tidens Løb og Historiens Gang opstod, var der en
betydningsfuld Aandsretning, der kun kunde tilfredsstille denne Trang ved at
døde den: Tidsløbet havde slet ingen Værdi, og det gjaldt kun om at stanse
det, at frelse sin Sjæl ud af det, gaa over i en tidløs Væren. Denne Vej gik
Udviklingen i Indien, især i Buddhismen, og den græske Tankegang, for
hvilken det Uforanderlige havde en større Værdi end det Foranderlige, gik i
beslægtet Retning. Den Tankegang derimod, som kom til at øve størst
Indflydelse paa den europæiske Menneskehed og paa hele den muhamedanske
Verden, anerkendte Tidsløbets Realitet, men fandt Frelsens Mulighed i en
stor Udviklingsgang, der fører Folk og Slægter til en sidste Afgørelse af
den vældige Verdenskamp mellem gode og onde Magter i og udenfor
Menneskeverdenen. Det synes at være fra den gammelpersiske Religion, at
denne Tanke især er udgaaet i Verden, hvad enten den er opstaaet indenfor
denne Religion eller gaar længere tilbage. Al Eschatologi, al Lære om en
Udvikling hen imod en yderste Dom og en derefter følgende Herlighedens eller
Fortabelsens Verden synes for en meget væsentlig Del at være af persisk
Oprindelse. Den har paavirket den senere Jødedom og derigennem
Kristendommen, og den skal (ifølge Goldzieher) direkte have paavirket
Muhamedanismen fra dennes første Begyndelse.

Ligesom Verdensdramaet, hvad Rummet angaar, fik sin begrænsede Skueplads i
det antike Verdensbillede, saaledes fik det, hvad Tiden angaar, sin klare og
bestemte Plads mellem Verdensbegyndelsen og Verdensafslutningen, og
Tidsafstanden mellem disse to Punkter, første Akts Begyndelse og sidste Akts
Slutning, var ligesaa lidt uendelig som den rumlige Afstand mellem Himmel og
Jord. Ogsaa her har Folkereligionen havt sin Styrke i Anskueligheden og den
med Anskueligheden følgende Begrænsning. Mennesket saa sig her med sit og
sine Fællers hele aandelige og materielle Ve og Vel som Led i en stor
Udviklingsgang, bestemt ved den fælles Fortid, maaske sukkende under
Nutidens Kval, men skuende med Fortrøstning hen til den Forløsning,
Fremtiden skulde bringe. Folkereligionen har her givet Livet en stor
Horizont, muliggjort en Forstaaelse af den Enkeltes Skæbne og Stræben fra et
højere Synspunkt. Og et betydningsfuldt Fremskridt i historisk Sans blev
derved gjort. Alt dette, uden at Anskuelighedens Ramme sprængtes. Det blev
muligt at danne klare Forestillinger om Fortid og Fremtid, Billeder, i
hvilke al den Spænding og Længsel saavel som al den Jubel og Sejrsfølelse,
Livet i Tiden kan udløse, kunde faa deres stærke, ofte ophøjede og
uforgængelige Udtryk.

Naar Eftertanken vaagner og gennemføres konsekvent, gaar det dog med denne
Anskuelighed, som med den, der raadede, hvad Rummet angik. De faste
Grænser, den absolute Begyndelse og den absolute Ende, falde, saasnart
Eftertanken anvender samme Metode paa Begyndelses- og Slutningsøjeblikkene
som paa de mellemliggende Øjeblikke. Istedetfor den korte Tid (ifølge
Parsismen 12,000 Aar), Verdensdramaet skulde vare, træder nu den
uoverskuelige Tid, baade hvad Fortid og Fremtid angaar. Og Spørgsmaalet
bliver da, hvorvidt der kan gives det Spand af Tid, Menneskets historiske
Bevidsthed kan omfatte, en saadan absolut Betydning, en saadan afrundet,
maalbestemt Karakter, som Folkereligionens Verdensopfattelse saa naturligt
frembød. Kan der nu overhovedet være Tale om et Verdensdrama? Den
Begrænsning, i hvilken den store Kunst viser sig, bliver umulig, hvad
Verdensløbet angaar, og dermed falder dettes dramatiske Karakter bort.

Ligesom ved Rumsbegrænsningen frembyder Anskueligheden ogsaa ved
Tidsbegrænsningen Muligheden af en Hjemlighedsfølelse. Det gaar jo saa med
Verden som med ethvert enkelt Væsen: den har sin Oprindelse, og den har sin
Død. At Verden var bleven gammel, var en almindelig Tro i de første kristne
Aarhundreder; ogsaa blandt Jøderne fandtes den, som det ses af det mærkelige
apokalyptiske Skrift „Fjerde Esra Bog“. Og man kunde da under Tidens Kvaler
og Skuffelser glæde sig til, at det stundede mod Enden. Den nære Afslutning
var da heller ingen Fantasi, som man fra Tid til anden svælgede i; den
prægede hele Livsanskuelsen, hele „det nye Testamentes Kristendom.“

Ligesom ved Rummet foregaar der dog ogsaa, hvad Tiden angaar, en
ejendommelig Inderliggørelsesproces i selve den religiøse Bevidsthed. Den
søger ud over Afstandene i Tid, som den søger ud over Afstandene i Rum. Ikke
i en fjern Fremtid eller Fortid, men i nærværende Oplevelse vil den have sin
Genstand. Idealet skal ikke være forsvundet eller blot fremtidigt, men gøre
sig gældende i selve det Liv, der leves i Nuet. Evigheden skal ikke ligge
før eller efter Tiden, men kunne erfares midt i Tiden, som jo ogsaa
psykologisk Erfaring godtgør, at Tiden forsvinder for os i de af store
Interesser og Ideer opfyldte Tilstande. I den religiøse Mystik, hvis Spirer
findes i selve de oprindelige religiøse Dokumenter, men som udvikledes
nærmere i forskellige Perioder af Religionens Historie, kommer den Tanke
atter og atter frem, at den sande Evighed kan opleves midt i Tiden, naar
Villien vender sig fra Alt, hvad der kan sprede, og samler sig om sin højeste
Genstand. Den religiøse Bevidsthed nærmer sig i saadanne Former til den
Opfattelse, der ene bliver mulig, naar Tidsbegrebet gennemtænkes.

Naar Eftertanken gennemføres, staa vi overfor den uoverskuelige Tid som
overfor det uoverskuelige Rum. Vi maa i begge Henseender finde vor Plads,
Pladsen for vort Arbejde og for de Erfaringer, vi kunne gøre om Livets Kaar.
Og ud fra denne Plads forfølge vi da Tidsløbet fremad og tilbage, saa langt
vi formaa. Men vi maa gøre os fortrolige med, at der ligger Tid og
Begivenheder før det Punkt, vi kunne begynde med, og efter det Punkt,
hvormed vi maa ende. Hvis vor Interesse ikke er begrænset til dette Spand af
Tid, der er forsvindende i Forhold til hvad den strenge Gennemførelse af
Tidsbegrebet vilde kræve, saa maa det gælde om at fastholde den
Overbevisning, at hvad der har virkelig Værdi for en overskuelig Tid, ikke
behøver at tabe den under det stadigt fortsatte Tidsløb, selvom denne Værdi
gennem mangehaande Forskydninger undergaar afgørende Ændringer. Indenfor de
Tidsrum, vi kunne overskue, se vi jo stadigt det Fænomen gentage sig, at
hvad der først kun stod som Middel, senere fik selvstændig Værdi, og at hvad
der først stod som et sidste Formaal, senere kom til at staa som Middel
eller Forberedelse for mere omfattende Værdier. Vi maa opgive at konstruere
Scenegangen i Verdensdramaet; vi kunne ikke afgøre, i hvilken Akt og i
hvilken Scene vort Liv falder; selve Ideen om et Drama kan kun anvendes paa
en begrænset Tid. Intet af de Dramer, Kunsten har frembragt, giver jo
iøvrigt ved sin Slutning nogen Sikkerhed for, hvorledes det gaar de
Overlevende senere hen. Og saaledes ere vi netop stillede overfor
Verdensløbet, naar vi lade Eftertanken raade. Videnskab er ikke Andet end
stadig og taalmodig Eftertanke.

\medskip

2. Med Udvidelsen af Verdensbilledet, hvad Rummet og Tiden angaar, fulgte i
den mærkelige Tidsalder, da den moderne Videnskab blev til, ogsaa Fordringen
om en strengere Gennemførelse af Aarsagsbegrebet, og dermed et nyt Ideal for
Erkendelse i Sammenligning med det, der havde raadet i Oldtid og
Middelalder. Den Grundsætning blev opstillet, at ligesom de Begivenheder
eller Forandringer, man søgte at finde Aarsagen til, forelaa som
Kendsgerninger i Erfaringen, saaledes skulde ogsaa de Aarsager, i hvilke man
fandt Forklaringen af dem, være saadanne, der kunde paavises i Erfaringen.
Kepler, af hvem denne Fordring især udtaltes med Eftertryk, havde selv i
sine tidligere Skrifter talt om Aarsager, der ikke fyldestgjorde denne
Fordring, og som han derfor senere ikke anerkendte som „\emph{gyldige
Aarsager}“ (veræ causæ). Han havde saaledes lært, at Planeterne holdtes og
førtes i deres Baner af Stjerneaander. Men han blev klar over, at denne
Forklaring videnskabeligt var ugyldig, da disse Aander ikke selv kunde
paavises i Erfaringen som virkende paa den Maade, han havde antaget. Han
ansaa derimod Magnetismen (som han sammenblandede med Tyngden) for en gyldig
Aarsag til, at Planeterne holdes i deres Baner, fordi den faktisk er en
tiltrækkende Kraft. Saa urigtig denne Opfattelse end var, saa var den dog i
den nye Forsknings Aand. Senere paaviste Newton, at den Kraft, der holder
Planeten i sin Bane, følger samme Lov som den Kraft, der lader den
udslyngede Sten falde ned paa Jorden. Det var et stort og forbilledligt
Exempel paa en „gyldig Aarsag“.

Ligesom ethvert Sted i Rummet maa bestemmes i Forhold til andre Steder, og
ethvert Øjeblik i Tiden i Forhold til andre Øjeblikke, saaledes maa i Kraft
af de gyldige Aarsagers Princip enhver Begivenhed bestemmes og forklares ved
dens Forhold til andre Begivenheder, der ere ligesaa paaviselige som den
selv. Herved kommer Erfaringens Verden til at staa som et sammenhængende
Hele, og Naturen fortolkes ved Naturen selv, ligesom man i en Bog fortolker
et Sted ved Hjælp af andre Steder i samme Bog. Vægten ligger atter her paa
Metoden, --- paa den Maade, paa hvilken man spørger, og den Kontrol, man
øver paa Svaret. Det er et nyt Organ, Menneskeaanden her har udviklet, og
som ethvert Organ kan det kun fungere paa sin bestemte Maade. Derved stilles
Opgaverne, og derved betinges Løsningerne. Og det er i Sammenhæng dermed en
ny intellektuel Trang, der har udviklet sig. Hidtil var det egentlig kun de
formale Videnskaber, Logik og Matematik, der havde udformet sig til
Selvstændighed; nu blev Erfaringsvidenskaben --- i de tre Hovedformer: som
Naturlære, Psykologi og Historie --- i Principet frigjort for teologiske og
spekulative Synsmaader. Og Udsigten over den nye Vej, der var aabnet, vakte
ofte en saadan Begejstring, at man paa dogmatisk Vis proklamerede, at Alt
kunde forklares efter de gyldige Aarsagers Princip, og at Aarsagsrækken var
absolut sammenhængende. Man gik Konstruktionens Vej istedetfor at vente, til
Iagttagelse og Forsøg kunde godtgøre, hvor Aarsagerne i de enkelte Tilfælde
laa. Det var en Forvexling af Resultat og Opgave, af Kendsgerninger og
Principer, af System og Metode. Men man tog, selv naar man kom ind paa
saadanne Overskridelser, ikke Fejl deri, at man var paa den rette Vej. Det
var Forskningens Aand, der var bleven sig sine Midler og sine Kræfter
bevidst og dermed tillige var bleven sig sit Ideal bevidst. Ti det var et
stort Ideal, der nu stillede sig for Erkendelsen: at forme et stedse
voxende Antal af Erfaringer til en stedse nøjere og inderligere
Enhedssammenhæng. Det er et Ideal, der til hver given Tid kun
tilnærmelsesvis kan fyldestgøres; det udraaber paa hvert Trin et Excelsior!
--- fører stedse ud til større Opgaver. Dogmatisk kan Videnskaben, naar den
forstaar sig selv, derfor ikke blive. Allerede Rummets og Tidens
Uoverskuelighed udelukker en Afslutning af vor Forsken. Der kan stedse
fremtræde noget Nyt. Og det ikke blot ud over Grænserne for det Rum og den
Tid, man til en given Tid har kunnet overskue, men ogsaa indenfor Grænserne.
Der kan opdukke nye Forskelle, hvor der hidtil synes at være Lighed eller
ubrudt Sammenhæng; en stedse rigere Fylde af Nuancer og Overgange kan der
opdages indenfor det, der syntes at være saa godt kendt. Verden kan stedse
blive større for os, baade i det Store og i det Smaa. Selve Begrebet Verden
kommer til at staa som et Idealbegreb, der ikke kan gennemføres, men stedse
aabner nye Horizonter. Og vi tør end ikke paastaa, at Tilværelsen er
begribelig eller forstaaelig, maalelig for vor Tanke. Den kunde være at
sammenligne med $\sqrt{2}$, det irrationelle Tal, der maaske kan bestemmes
fremskridende gennem Fund af stedse flere Decimaler, men ikke kan faa nogen
afsluttende Bestemmelse. Begrebet Verden som afsluttet Hele er egentlig slet
ikke noget videnskabeligt Begreb. Det er et Begreb, den religiøse Bevidsthed
har dannet i sin Utaalmodighed, for at faa en skarp Modsætning til sin
Gudsforestilling, men som for Videnskaben ikke betegner noget Givet, kun en
stadig Opgave. For Videnskaben er der ligesaa stor Vanskelighed ved Begrebet
Verden som ved Begrebet Gud.

Naar Religionen i sine klassiske Tider ikke blot gav Hjælp og Trøst under
Livets Gang, men ogsaa tilfredsstillede Erkendelsestrangen, saa var det ikke
de „gyldige Aarsagers“ Princip den opfordrede til at anvende. Dens
Naturforklaring var \emph{animistisk}, d.\ v.\ s. den hævdede personlige
Væseners Indgriben som Forklaringsgrund. Ligesom Verden som Helhed antoges
for skabt ved en Villiesakt, saaledes antoges ogsaa de enkelte Begivenheder
i Naturen, i Sjælelivet og i Historien bevirkede ved personlige Væseners
Handlen. Især de i Menneskeskæbner stærkt indgribende Naturbegivenheder
forklaredes paa denne Maade. Uvejr og Jordskælv saavel som Regn og frugtbare
Tider fandt deres tilstrækkelige Grund i Forsynets Villie, der snart
straffede eller prøvede, snart lønnede og velsignede. Afgørende Forandringer
paa Sjælelivets Omraade forklaredes ikke ved psykologiske Aarsager, men ved
guddommelig Indgriben, som naar Aarsagen til, at Achilleus betvinger sin
Vrede overfor Agamemnon, forklares ved, at Atene sendes ned fra Olymp for at
tale ham til Rette. Historiens Gang tydes af Profeterne som Udførelse af en
guddommelig Plan, der bøjer Fyrsters Hjerter og vender Folkenes Skæbner efter
sine Formaal. I afbleget Form optræder den animistiske Forklaring i den
saakaldte teleologiske Forklaringsmaade, i hvilken man paa mere ubestemt
Maade taler om Formaalsaarsager, uden at være saa konsekvent at betragte
disse Formaal som Udtryk for personlige Villier.

Næppe nogensinde er dog den animistiske Naturforklaring bleven konsekvent
gennemført. Allerede daglig Erfaring ledte til at antage en fast Sammenhæng
mellem Begivenheder, der begge forelaa som Kendsgerninger, saa at den ene
kom til at staa som Aarsag til den anden. Man lod den begyndende Forsken gaa
ad sine egne Veje et længere eller kortere Stykke, men betrag-
% [transcription continues: printed p. 21 / PDF p. 30 — mid-word "betrag-tede…"]

\end{document}
